\scanpage{41}
% archive foliation: 40
\[
\operatorname{Tor}_p^B(M_1,-M_n)\otimes_A A'-\operatorname{Tor}_p^{B'}(M'_1,-,M'_n)
\]
est un isomorphisme pour $p\leq\unsure{n|N}$. Or c'est précisément l'hom considéré dans la proposition.

\uline{Corollaire} \inserted{Sous les conditions de la prop.} Il existe une décomposition cellulaire de $S$ en des $S_\alpha$ loc. fermés, ($1\leq i\leq n$) tels que, posant \inserted{$Y^{\alpha}=Y\times_S S_\alpha$} $X_i^{(\alpha)}=X_i\times_S S_\alpha$, $F^{(\alpha)}=F\otimes_S S_\alpha$, \deleted{\unsure{on ait pour tt}} qu'on ait ceci : pour tt changement de base $T-S_\alpha$, l'hom canonique
\[
\operatorname{Tor}_p^{\underline O_{S_\alpha}}(F_1^{(\alpha)},-,F_n^{(\alpha)})\otimes_{S_\alpha}T\longrightarrow\operatorname{Tor}_p^{\underline O_T}(F_{1(T)}^{(\alpha)},-,F_{n(T)}^{(\alpha)})
\]
est un isom pour $|p|\leq N$.

\margin{\figure{\begin{tikzpicture}\draw (0,0)--(.13,-3);\draw (.2,-4.4)--(.24,-4.75);\end{tikzpicture}}}
\uline{Proposition} $X$ de type fini sur $S$ noeth, $F$ et $G$ deux faisceaux \inserted{coh} sur $X$. \inserted{$N$ un entier.} Il existe alors une partie ouverte non vide $U$ de $S$ telle que, posant $S'=U_{\text{réd}}$, $X'=X\times_S S'$, $F'=F\times_{\unsure{S}}S'$, $G'=G\times_S S'$, l'on ait ceci : pour tt changement de base \inserted{$T\longrightarrow S$}, on a un isomorphisme
\[
\underline{\operatorname{Ext}}^p_{\underline O_{S'}}(F',G')\otimes_{S'}T\longrightarrow\underline{\operatorname{Ext}}^p_{\underline O_T}(F'_T,G'_T).
\]
pour $|p|\leq\unsure{n|N}$
